\section{Constructing one marked arithmetic frame and its two strict classes}
\label{sec:marked-frame-construction}

\paragraph{The result and its precise boundary.}
This chapter constructs the actual auxiliary field, marked element,
profinite quotient, local groups and two nonzero strict cohomology
classes used by \texttt{frame\_ex}.  It proves the complete source
suppliers \texttt{exists\_\allowbreak twistCharacter\_\allowbreak continuousCohomology},
\texttt{hlay}, \texttt{l15}, \texttt{ufacts}, \texttt{b0},
\texttt{link\_ex}, \texttt{ok\_basic}, \texttt{strict} and
\texttt{frobenius\_\allowbreak galoisField\_\allowbreak four\_\allowbreak compatibility}.
The construction uses one field and one marked element throughout.
The real quadratic auxiliary field admits a useful elementary
simplification: a Chinese-remainder construction replaces the
source's prime in a prescribed progression.

The entry premise is the \emph{full} original \texttt{Gen} record,
not just its seed-relevant subset from CAF.  The strict-class step
retains exactly the original \(P:\texttt{GC.Hypotheses}_{2}\).
Its finite local and global cohomology laws and Greenberg--Wiles
identity are the laws used below.  The uniform class-rank step
uses the already reconstructed TW1 and its named classical inputs.
No automorphic rank-two or rank-three seed is assumed or produced
by this construction.

At the end, the resulting frame has \texttt{Big}, both strict
classes, and twenty-one of the twenty-eight literal \texttt{Ok}
fields.  The seven remaining fields are named explicitly at the
end of the chapter.  Thus the chapter produces several complete
upstream theorems and an actual shared frame, but does not yet
claim the complete \texttt{HasFrameW}, \texttt{GdFrameEx},
\texttt{HostFrameEx} or \texttt{GdEnd} endpoints.

\subsection{Notation and the exact arithmetic input}

Let \(\Omega\) be the specified algebraic closure of \(\mathbb Q\)
and \(G=\operatorname{Gal}(\Omega/\mathbb Q)\).  Fix an odd prime
\(\ell\), a commutative local ring \(O\) with finite residue field
\(k\) of characteristic \(\ell\), and characters
\[
 \eta:G\longrightarrow\{\pm1\},\qquad
 \Psi:G\longrightarrow O^\times,\qquad
 \alpha=\overline\Psi:G\longrightarrow k^\times.
\]
Write \(\epsilon:G\to\mathbb Z_\ell^\times\) for the actual
cyclotomic character, \(\epsilon_r\) for its reduction modulo
\(\ell^r\), and \(\omega\) for its reduction in \(k^\times\).
The prime-field map is injective, so
\(\ker\omega=\ker\epsilon_1\).
For a valuation subring \(A\subset\Omega\), write \(D_A,I_A\)
for its decomposition and inertia groups, and write \(A\mid n\)
for \(A(n)<1\), with the source's multiplicative valuation.

We fix \(\sigma,B,M,H_b\) satisfying the original \texttt{Gen}:
\(M>0\), \(2\ell\mid M\), \(\sigma\) is a complex conjugation,
\(B\mid\ell\), and there are \(D>0,\delta\in\Omega\) with
\[
 \delta^2=-D,\qquad s\delta=\eta(s)\delta,
 \qquad p\mid D\Longrightarrow p\mid M
 \quad\text{for every rational prime }p.                  \tag{MFC.1}
\]
The characters \(\eta,\alpha\) factor through the action on
\(\mu_{M^a}\) for some \(a\); \(\eta\) is unramified at
\(\ell\); \(\alpha(\sigma)=1\); and one integer \(j\ge0\)
satisfies \(\alpha|_{I_A}=\omega^j|_{I_A}\) for every
\(A\mid\ell\).  The two separation hypotheses are
\[
 \ker\omega\cap\ker\alpha\not\subseteq\ker\eta,
 \qquad
 \exists t\in D_B\cap\ker\eta:\ \alpha(t)^2\ne1.          \tag{MFC.2}
\]
Finally \(H_b\) is open, contains \([G,G]\), has index
\(\ell^b\) for some \(b\ge0\), contains every \(I_A\) with
\(A\mid\ell\), and satisfies the exact lifting condition
\[
 A\nmid\ell,\quad s\in I_A\cap H_b,\quad\alpha(s)=1
                  \quad\Longrightarrow\quad\Psi(s)=1.      \tag{MFC.3}
\]
These are all nineteen source fields, with the finite residue-field
and odd-prime conditions included.  In particular (MFC.2)'s second
condition and (MFC.3) are not inferred from the weaker CAF entry
conditions.

The subgroup \(H_\theta=H_\theta(w)\) is the inverse image of
the \(\ell\)-th powers in \((\mathbb Z/w\mathbb Z)^\times\).
CAF4 proved the literal source \texttt{Hth} identification and its
open, abelian, index-\(\ell\), unramified and inertia properties.
A witness prime \(w\) means exactly
\[
 w\text{ prime},\quad w\nmid M,\quad w\equiv1\pmod\ell,
 \quad D_A\subseteq\ker\eta\cap\ker\alpha\cap\ker\omega
                     \quad(A\mid w).                       \tag{MFC.4}
\]
CAF5 supplies such a prime with \(w\equiv1\pmod Q\), for every
already prescribed positive \(Q\).

\subsection{The real quadratic subgroup and the entire tower}

% BEGIN MFC1
\paragraph{MFC1. An elementary real quadratic auxiliary subgroup.}
\label{claim:MFC1}
There is an open subgroup \(H_0\subset G\) of index two such that
every complex conjugation belongs to \(H_0\),
\[
 D_A\cap H_0=D_A\cap\ker\eta\quad(A\mid M),                \tag{MFC.5}
\]
and some valuation \(A\nmid M\) has \(I_A\not\subseteq H_0\).
These are all fields of \texttt{FrobeniusSplitPrimeDatum}.

\begin{proof}
Choose an odd prime \(p\nmid2DM\).  Such a prime exists by
Euclid's argument.  The Chinese remainder theorem gives a positive
integer \(n\) with
\[
                  n\equiv-1\pmod{4M},\qquad
                  n\equiv p\pmod{p^2}.                    \tag{MFC.6}
\]
Choose \(x\in\Omega\) with \(x^2=-n\), put \(z=\delta x\),
and let \(H_0\) be the stabilizer of \(z\).  Since
\(z^2=Dn>0\) and \(v_p(Dn)=1\), \(Dn\) is not a rational
square.  Hence \(\mathbb Q(z)/\mathbb Q\) is quadratic and
\(H_0\) is open of index two.  Every complex embedding sends
\(z\) to one of the real roots of \(X^2-Dn\), so every complex
conjugation fixes \(z\).

Write \(-n=1+4Mt\) with \(t\in\mathbb Z\).  We claim that
every \(s\in D_A\), with \(A\mid M\), fixes \(x\).
It sends \(x\) to \(x\) or \(-x\).  In the latter case put
\(y=(1+x)/2\).  Then \(sy=1-y\) and
\(y(1-y)=-Mt\).  Decomposition-group invariance makes the
valuations of \(y\) and \(1-y\) equal, while their product
has valuation less than one because integers have valuation at
most one.  Both valuations are therefore less than one.
Their sum is one, contradicting the nonarchimedean inequality.
Thus \(sx=x\).  Since \(z\ne0\), (MFC.1) now gives
\(sz=z\) if and only if \(\eta(s)=1\), proving (MFC.5).

At the prime \(p\), the polynomial \(X^2-Dn\) is Eisenstein.
The quadratic extension is totally ramified there, so its inertia
group has a nonidentity image and sends \(z\) to \(-z\).
For any corresponding extension \(A\) of the \(p\)-adic
valuation, \(I_A\not\subseteq H_0\) and \(A\nmid M\).
This proves the final field.  No requirement that the auxiliary
integer \(n\) itself be prime was needed.
\end{proof}
% END MFC1

% BEGIN MFC2
\paragraph{MFC2. Exact cyclotomic layers and their local indices.}
\label{claim:MFC2}
For \(m\ge0\), define the literal source subgroup
\[
 H_m=\ker(\epsilon_{m+1}^{\ell-1}).                         \tag{MFC.7}
\]
It is open, contains \([G,G]\), has index \(\ell^m\), and
\[
 H_m\cap\ker\epsilon_1=\ker\epsilon_{m+1}.                \tag{MFC.8}
\]
It contains inertia at every valuation not above \(\ell\).
For every \(A\mid\ell\),
\[
 [I_A:I_A\cap H_m]=\ell^m,
 \qquad\epsilon_r(I_A)=(\mathbb Z/\ell^r\mathbb Z)^\times
                         \quad(r\ge0).                    \tag{MFC.9}
\]
These are the full conclusions of \texttt{hlay}.

\begin{proof}
For \(r\ge1\), \(\Phi_{\ell^r}(1+T)\) is Eisenstein at
\(\ell\): its reduction is a power of \(T\) and its constant
term is \(\ell\).  Thus \(\mathbb Q_\ell(\mu_{\ell^r})\)
is totally ramified of its full degree.  The inertia action on
these roots is consequently the whole unit group, at every
prolongation of the valuation.  The assertion at \(r=0\) is
trivial.  This also proves global surjectivity of \(\epsilon_r\).

For odd \(\ell\), the units modulo \(\ell^{m+1}\) have a
subgroup of order \(\ell-1\) consisting of the unique lifts
of the nonzero residues satisfying \(u^{\ell-1}=1\).
Existence and uniqueness follow from Hensel's lemma, since the
derivative \((\ell-1)u^{\ell-2}\) is a unit.  Its complement
is the subgroup \(1+\ell\mathbb Z/\ell^{m+1}\mathbb Z\),
of order \(\ell^m\).  Raising to \(\ell-1\) kills the first
factor and is an automorphism of the second.  Hence its image
has order \(\ell^m\), proving the index assertion and (MFC.9).
If a unit is also one modulo \(\ell\), its order is an
\(\ell\)-power.  If its \((\ell-1)\)-st power is one, its
order is one.  This proves (MFC.8).

The kernel of the finite cyclotomic action is open, so \(H_m\)
is open.  Its quotient is abelian.  Away from \(\ell\), all
\(\ell\)-power roots are unramified: their defining polynomials
have invertible derivative over the residue field, and Hensel
lifts their roots from an unramified extension.  Hence such
inertia is contained even in \(\ker\epsilon\), as asserted.
\end{proof}
% END MFC2

% BEGIN MFC3
\paragraph{MFC3. Residual-character identities cannot appear only in the tower.}
\label{claim:MFC3}
Fix \(H_0\) as in MFC1 and a witness prime \(w\).  Put
\[
 U_m=H_0\cap H_\theta(w)\cap H_b\cap H_m.                 \tag{MFC.10}
\]
For every \(i,j\in\mathbb Z\), the character
\(\chi=\omega^i\alpha^j\) has the following properties:
an identity \(\chi=1\) on \(U_m\) holds on \(G\);
an identity \(\chi=\eta_k\) on \(U_m\) holds on \(G\);
and an identity \(\chi=1\) on \(U_m\cap\ker\eta\) holds
on \(\ker\eta\).  Here \(\eta_k\) is the same quadratic
character viewed in \(k^\times\).  This is the full
\texttt{l15} conclusion.

\begin{proof}
Let \([G:H_b]=\ell^b\).  Each of \(H_\theta,H_b,H_m\)
is normal, with abelian quotient of order respectively
\(\ell,\ell^b,\ell^m\).  Lagrange's theorem gives
\[
 s^{\ell^{1+b+m}}\in U_m\qquad(s\in H_0).                 \tag{MFC.11}
\]
An element of \(k^\times\) whose \(\ell^N\)-th power is one
is one: \((u-1)^{\ell^N}=u^{\ell^N}-1=0\) in a field.
Therefore a character trivial on \(U_m\) is trivial on
\(H_0\), and the same argument applies on intersections with
any subgroup.

Choose \(A\nmid M\) and \(t\in I_A\setminus H_0\) from
MFC1.  The cyclotomic conductor hypotheses imply
\(\eta(t)=\alpha(t)=1\): inertia away from \(M\) fixes
\(\mu_{M^a}\).  Since \(\ell\mid M\), also
\(\omega(t)=1\).  Thus \(\chi(t)=\eta_k(t)=1\).
Every element outside the index-two subgroup \(H_0\) is the
product of an element in \(H_0\) with \(t\), or its inverse.
This extends the first identity to all of \(G\).  Apply the
same argument to \(\chi\eta_k^{-1}\) for the second identity.
For the third, note that \(t\in\ker\eta\) and apply the
identical coset argument inside \(\ker\eta\).
\end{proof}
% END MFC3

% BEGIN MFC4
\paragraph{MFC4. All twelve fields of the original \texttt{UFacts}.}
\label{claim:MFC4}
For every \(m\ge0\), \(U=U_m\) satisfies the following
complete source record:
\begin{enumerate}
\item \(U\) is open and contains \([G,G]\).
\item Every complex conjugation lies in \(U\).
\item \(I_A\cap U\subseteq\ker\eta\) for every valuation \(A\).
\item \(D_A\cap U\subseteq\ker\eta\) whenever \(A\mid M\).
\item If \(A\nmid\ell\), \(s\in I_A\cap U\), and
      \(\alpha(s)=1\), then \(\Psi(s)=1\).
\item If \(A\mid\ell\), some \(s\in I_A\cap U\) has
      \(\epsilon_1(s)\ne1\).
\item If \(A\mid\ell\), \([D_A:D_A\cap U]\ge\ell^m\).
\item For \(j=1,-1\), some \(s\in U\cap\ker\eta\) has
      \(\omega(s)\alpha(s)^j\ne1\).
\item The same assertion holds for \(j=2,-2\).
\item At every \(A\mid\ell\), some \(s\in D_A\cap U\)
      has \(\alpha(s)^2\ne1\).
\item There is no \(j\in\mathbb Z\) with
      \(\omega(s)=\alpha(s)^j\) for every
      \(s\in U\cap\ker\eta\).
\end{enumerate}
The first item contains the two separate fields
\texttt{isOpen} and \texttt{ab}; the remaining items are
\texttt{real}, \texttt{unr}, \texttt{split}, \texttt{teich},
\texttt{mu}, \texttt{deg}, \texttt{f2b}, \texttt{f2c},
\texttt{f3} and \texttt{f5}.

\begin{proof}
Intersections of finitely many open subgroups are open.
Every index-two subgroup contains the commutator, as do the other
three factors in (MFC.10).  A complex conjugation has order two.
Its image in a group of odd order is trivial.  Hence it lies in
the three \(\ell\)-power-index factors and, by MFC1, in \(H_0\).

At a valuation above \(M\), (MFC.5) gives both the inertia and
decomposition assertions.  At every other valuation, \(\eta\)
is already trivial on inertia by its cyclotomic conductor.
The lifting assertion follows from \(U\subseteq H_b\) and
(MFC.3).

For \(A\mid\ell\), choose \(s\in I_A\) with
\(\epsilon_{m+1}(s)=-1\), using MFC2.  It belongs to \(H_m\)
because \(\ell-1\) is even.  Since \(\eta\) is unramified
at \(\ell\), (MFC.5) puts it in \(H_0\); here \(A\mid M\)
because \(\ell\mid M\).  It lies in \(H_b\) by hypothesis
and in \(H_\theta\) because \(w\ne\ell\).  Thus \(s\in U\)
and \(\epsilon_1(s)=-1\ne1\).  The index estimate follows
from \(I_A\subseteq D_A\), \(U\subseteq H_m\), and
\([I_A:I_A\cap H_m]=\ell^m\).

By (MFC.1), \(\eta(\sigma)=-1\), since a complex conjugation
negates a nonzero square root of a negative rational number.
Choose \(t\in\ker\omega\cap\ker\alpha\) with
\(\eta(t)=-1\) from (MFC.2).  Then
\[
 s_*=t\sigma\in\ker\eta,\qquad
                  \omega(s_*)=-1,\quad\alpha(s_*)=1.       \tag{MFC.12}
\]
If \(\omega\alpha^j=1\) throughout \(U\cap\ker\eta\),
MFC3 extends this identity to \(\ker\eta\), contradicting
(MFC.12).  In fact this proves the stated \texttt{f2b} and
\texttt{f2c} assertions for every integer \(j\).  Applying it
to \(\omega\alpha^{-j}\) also proves \texttt{f5}.

Finally take the element in the second assertion of (MFC.2).
Conjugate it so that its decomposition group is \(D_A\) for
the requested \(A\mid\ell\).  Characters into abelian groups
are invariant under conjugation, so it still has \(\eta=1\)
and \(\alpha^2\ne1\).  Equation (MFC.5) puts it in \(H_0\).
Raise it to \(\ell^{1+b+m}\).  Equation (MFC.11) puts the
result in \(U\), without leaving \(D_A\).  Its residual
square is still nonidentity by the characteristic-\(\ell\)
argument following (MFC.11).  This proves \texttt{f3} at every
prolongation, completing the record.
\end{proof}
% END MFC4

% BEGIN MFC5
\paragraph{MFC5. One class-rank bound before the layer is chosen.}
\label{claim:MFC5}
For the fixed \(H_0,w,H_b\), there is \(C_0\in\mathbb N\)
such that, for every \(m\ge0\),
\[
 \operatorname{clRk}_\ell
  (U_m\cap\ker\eta\cap\ker\alpha\cap\ker\epsilon_1)
                         \le C_0.                         \tag{MFC.13}
\]
Here \(\operatorname{clRk}\) is exactly \texttt{GC.clRk}:
the \(\mathbb F_\ell\)-rank of the locally constant additive
characters of the indicated subgroup which vanish on every
intersection with inertia.  This is the full \texttt{b0}
supplier, with \(C_0\) chosen before \(m\).

\begin{proof}
Set
\[
 V=H_0\cap H_\theta\cap H_b\cap\ker\eta\cap\ker\alpha
                                      \cap\ker\epsilon_1.
\]
This subgroup is open: the two character kernels contain the
kernel of a finite cyclotomic action.  It contains \([G,G]\).
Thus \(E=\Omega^V\) is a finite abelian number field, and it
contains \(\mu_\ell\).  By (MFC.8),
\[
 U_m\cap\ker\eta\cap\ker\alpha\cap\ker\epsilon_1
                 =V\cap\ker\epsilon_{m+1}
                 =G_{E(\mu_{\ell^{m+1}})}.                \tag{MFC.14}
\]
These are equal subgroups of the original \(G\), not merely
abstractly isomorphic Galois groups.

The final decomposition condition in the source's
\(\texttt{Kum.XDual}\ 1\) is vacuous because no valuation
has \(A(1)<1\).  Its inertia intersections are the inertia
groups for the actual fixed field in (MFC.14).  Finite class-field
reciprocity therefore identifies its additive character space with
\[
 \operatorname{Hom}(\operatorname{Cl}(E(\mu_{\ell^{m+1}})),
                    \mathbb F_\ell).
\]
There is no real-place qualification: \(E\) contains
\(\mu_\ell\) for odd \(\ell\), hence is totally complex.
This Hom space has dimension
\(\log_\ell|\operatorname{Cl}(E(\mu_{\ell^{m+1}}))/\ell|\),
by elementary finite abelian group theory.  TW1 gives one bound
\(\ell^{C_0}\) for these class-group quotients at all levels,
including repetitions at the beginning of the tower.  This proves
(MFC.13) at its exact fixed-field and coefficient scope.
\end{proof}
% END MFC5

% BEGIN MFC6
\paragraph{MFC6. The field and the marked local element are chosen together.}
\label{claim:MFC6}
After MFC5 choose \(m=3C_0+7\), put \(U=U_m\), and take
the actual field
\[
                  L=\Omega^{U\cap\ker\eta},\qquad F=L^+.
                                                               \tag{MFC.15}
\]
Then \(L/\mathbb Q\) is finite abelian CM,
\(G_F=U\), \(G_L=U\cap\ker\eta\), and there is
\[
              g_1\in D_B\cap G_L,\qquad\alpha(g_1)^2\ne1.
                                                               \tag{MFC.16}
\]
Moreover \(\ell^m>3C_0+7\).

\begin{proof}
MFC4 supplies precisely the \texttt{UFacts} input to CAF8.
That theorem gives (MFC.15), including its literal Galois-group
equalities and maximal real subfield.  Use MFC4's \texttt{f3}
at the original \(B\) to choose \(g_1\in D_B\cap U\) with
the required residual square.  Since \(B\mid M\), the
\texttt{split} field gives \(\eta(g_1)=1\), proving (MFC.16).
Finally \(m\ge7\), and \(\ell^m>m\) follows by induction
from \(\ell\ge3\).  No constant has been chosen after its
own layer: the order is \(H_0,w,C_0,m,L,g_1\).
\end{proof}
% END MFC6

\subsection{The literal global quotient and local frame}

% BEGIN MFC7
\paragraph{MFC7. All local places, embeddings and degrees on the same field.}
\label{claim:MFC7}
Let \(L,F,U,g_1\) satisfy MFC6; the same assertions below hold
for any original \(L\) and \(g_1\) satisfying the hypotheses
of \texttt{link\_ex}.  Let \(\mathcal V\) be the set of primes
of \(F\) above \(\ell\).  One may choose prolongations
\(A_v\) and elements \(\gamma_v\in\ker\eta\) with
\[
 A_v=\gamma_v B,\qquad
 G_v=D_{A_v}\cap U,\qquad I_v=I_{A_v}\cap U.              \tag{MFC.17}
\]
The set \(\mathcal V\) is finite and nonempty.  Every valuation
above \(\ell\) is \(U\)-conjugate to exactly one \(A_v\),
and, with the natural embeddings \(j_v:G_v\hookrightarrow U\),
\[
 d_v=[F_v:\mathbb Q_\ell]=[D_{A_v}:D_{A_v}\cap U],
 \qquad \sum_vd_v=[F:\mathbb Q].                           \tag{MFC.18}
\]
There are also complex conjugations \(c_\nu\in G_F\) at
every real place \(\nu\) of \(F\).

\begin{proof}
Primes above a fixed nonzero rational prime are finite and nonempty
in a number field.  Choose an extension \(A\) of each.  All
prolongations of the rational \(\ell\)-adic valuation to the
normal algebraic closure are conjugate, so \(A=\gamma B\)
for some \(\gamma\in G\).  If \(\eta(\gamma)=-1\), replace
\(\gamma\) by \(\sigma\gamma\) and \(A\) by \(\sigma A\).
Here \(\sigma\in U=G_F\), so the prime of \(F\) is unchanged,
and \(\eta(\sigma\gamma)=1\).  This gives the required choices.

Prolongations of one valuation of \(F\) are conjugate by \(G_F\).
Conversely a member of \(G_F\) cannot change the restricted
valuation on \(F\).  This proves both existence and uniqueness
of the index \(v\) in the asserted orbit description.  Restriction
identifies the decomposition and inertia groups over \(F\) with
the intersections in (MFC.17).  Completion of the finite extension
\(F/\mathbb Q\) identifies its local Galois degree with the
decomposition quotient in (MFC.18).  Summing the local degrees
gives the global degree, for example by the decomposition of
\(F\otimes_\mathbb Q\mathbb Q_\ell\) into its finite local
factors.

For a real embedding \(F\hookrightarrow\mathbb R\), extend
it to \(\Omega\hookrightarrow\mathbb C\).  Conjugation on
the image preserves the algebraic closure and fixes \(F\),
so it pulls back to the required \(c_\nu\in G_F\).  This
also records the exact real place at which it is a conjugation.
\end{proof}
% END MFC7

% BEGIN MFC8
\paragraph{MFC8. The complete original \texttt{Linked} frame.}
\label{claim:MFC8}
Under the full hypotheses of \texttt{link\_ex}, there is an
actual \(\Phi:\texttt{Frame}\ \ell\ k\) satisfying
\(\texttt{Linked}\ \ell\ \eta\ \Psi\ \sigma\ g_1\ B\ M\ L\ \Phi\).
In particular it contains every \texttt{GalLink} and
\texttt{GC.FLink} field, with the same marked \(g_1\), field,
characters and local maps.

\begin{proof}
Identify \(G_F\) with \(U=\operatorname{upl}(\Omega,L)\)
by restriction to the same \(\Omega\); this is a topological
group isomorphism.  Its inverse merely regards an automorphism
fixing the embedded \(F\) as an \(F\)-automorphism.  Define
\[
 K_0=U\cap\ker\eta\cap\ker\alpha,
 \qquad
 J=\overline{\langle\!\langle
       I_A\cap K_0:\ A\nmid\ell
                   \rangle\!\rangle}_{,U},\qquad
                         \Gamma=U/J.                      \tag{MFC.19}
\]
The normal closure and closure are taken in this particular \(U\).
In the \(G_F\) notation, this is exactly
\(\texttt{tameN}(S,K_0)\), with \(S\) the full finite set
of primes above \(\ell\).  Every nontrivial valuation restricts
to a finite prime of \(F\); the trivial valuation has trivial
inertia and adds only the identity to the generating set.
Thus the two descriptions have equal normal closures, including
the valuation quantifier in \texttt{GalLink.p\_ker}.

The subgroup \(K_0\) is open and normal in \(U\).  Its index there divides
\(2|k^\times|\), which is prime to \(\ell\), since its
quotient is the image of \((\eta,\alpha)\).  Because \(K_0\)
is closed and normal, \(J\subseteq K_0\).  Also
\(J\subseteq\ker\epsilon\), since \(\epsilon\) is
unramified away from \(\ell\).  Consequently \(\eta,\alpha\)
and \(\epsilon\) descend to characters
\[
 \delta_\Phi:\Gamma\to\{\pm1\},\quad
 \psi_\Phi:\Gamma\to k^\times,\quad
 \epsilon_\Phi:\Gamma\to\mathbb Z_\ell^\times.            \tag{MFC.20}
\]
Let \(p:U\twoheadrightarrow\Gamma\) be the canonical quotient.
It is continuous and onto.  The quotient is profinite: a closed
normal subgroup of a profinite group is the intersection of its
open normal overgroups, and the quotient is their inverse limit
of finite quotients.  This supplies its compactness, Hausdorff
property and total disconnectedness.

Use MFC7's local groups, with \(\operatorname{loc}_v=pj_v\),
the indicated \(I_v\), degrees \(d_v\), real-place index set,
and \(c_w=p(c_\nu)\).  Since \(\ell\mid M\), \texttt{UFacts.split}
says \(\eta=1\) on each \(G_v\).  Hence the source's
\(\operatorname{loc0}_v:G_v\to\ker\delta_\Phi\) is well
defined.  Set
\[
                    c=p(\sigma),\qquad g_0=p(g_1).
                                                               \tag{MFC.21}
\]
The latter is regarded as an element of \(\ker\delta_\Phi\),
since \(g_1\in G_L\).  Together (MFC.17)--(MFC.21) specify
every field of the original \texttt{Frame} record.

We verify the full linking records, to retain the actual maps.
The field is abelian and \(G_L=U\cap\ker\eta\) by hypothesis.
The \texttt{unr}, \texttt{split} and \texttt{teich} fields are
the corresponding \texttt{UFacts} fields.  The map \(p\), its
continuity, surjectivity and exact kernel are (MFC.19).  Its three
character identities are precisely
\[
 \delta_\Phi(pu)=\eta(u),\quad
 \epsilon_\Phi(pu)=\epsilon(u),\quad
 \psi_\Phi(pu)=\alpha(u).                                 \tag{MFC.22}
\]
The \texttt{c} and \texttt{g0} equations are (MFC.21).
The chosen \(A_v,\gamma_v\) give \texttt{pl}, \texttt{pl\_conj}
and \texttt{pl\_all}.  Inclusion gives continuous injective
\(j_v\), its exact decomposition-group range, the equality
\(\operatorname{loc}_v=pj_v\), and the exact inertia preimage.
Equation (MFC.18) gives \texttt{dv}.  Finally
\texttt{UFacts.mu} produces an element in \(I_{A_v}\cap U\)
whose residual cyclotomic character is nontrivial.  Viewed in
\(G_v\), it gives \texttt{GalLink.mu}.

For \texttt{FLink}, take \(S=\mathcal V\), the original
\(G_F\)-version of \(K_0\), the canonical quotient map
\(\pi:G_F\to\Gamma\), and the decomposition inclusions
\(\iota_v:G_v\hookrightarrow G_F\).  Take the prime map
\(v\mapsto v\), the identity real-place equivalence, and the
actual \(c_\nu\) from MFC7.  The \texttt{GLink} fields are:
odd \(\ell\), all primes above \(\ell\), normal open
\(K_0\) of prime-to-\(\ell\) index, continuous surjective
\(\pi\), and kernel exactly \texttt{tameN}.  All have just
been verified.  The \texttt{LLink} fields are MFC7's place,
range and inclusion statements.  Its remaining equations are
the definitions of \(I_v,d_v,\epsilon_\Phi,c_w\).
The quadratic kernel is open because its inverse image under
the quotient map is open; \(\delta_\Phi(c)=-1\) follows
from \(\eta(\sigma)=-1\).  Finally both \(\delta_\Phi\)
and \(\psi_\Phi\) kill \(\pi(K_0)\).
The two final comparison equations in \texttt{Linked} say
that \(\pi\) and \(p\), and \(\iota_v\) and \(j_v\),
are these same maps under the fixed topological identification
\(G_F=U\).  They hold by construction.  This proves the full
source endpoint without an additional linking assumption.
\end{proof}
% END MFC8

% BEGIN MFC9
\paragraph{MFC9. All original basic frame conclusions and the size bound.}
\label{claim:MFC9}
For any linked frame in MFC8 and any marked
\(g_1\) with \(\alpha(g_1)^2\ne1\), all conclusions of
\texttt{ok\_basic} hold.  They include the nineteen
\texttt{Ok} fields other than \texttt{f4}, \texttt{f4'}
and the seven generic cohomological fields listed below.
They also give \(d_v\ge\ell^m\) for every \(v\) and
\(\mathcal V\ne\varnothing\).  With MFC6's \(m\),
\[
       3C_0+7<d=\sum_vd_v,\qquad5<d_v\quad\text{for all }v,
                                                               \tag{MFC.23}
\]
which is exactly \texttt{Big}\(\ C_0\).

\begin{proof}
For an arbitrary linked frame, each equality used here is one
of its \texttt{GalLink} or \texttt{FLink} fields; the argument
does not require that its underlying type be the particular
quotient model constructed in MFC8.
The characters \(\delta_\Phi,\psi_\Phi\) have open kernel
because their pullbacks in (MFC.22) have open kernel and \(p\)
is a quotient map.  The same quotient property gives continuity
of \(\epsilon_\Phi\); the local maps are continuous by their
definitions.  Complex conjugation has square one, giving
\(c^2=1\), \(c_w^2=1\).  Its action on \(\delta\) gives
\(\delta_\Phi(c)=\delta_\Phi(c_w)=-1\), and the original
evenness gives \(\psi_\Phi(c)=1\).  The marked equation is
\(\psi_\Phi(g_0)^2=\alpha(g_1)^2\ne1\).
Oddness of \(\ell\) and finiteness of \(k\) are the input
\texttt{Gen} fields.  Since \(F\) is totally real, MFC7 gives
\(\#\operatorname{Inf}=[F:\mathbb Q]=\sum_vd_v\).

Reduction of (MFC.22)'s cyclotomic identity gives
\(\operatorname{om}_\Phi(pu)=\omega(u)\).  This is the
source's exact residual character: the two maps
\(\mathbb Z_\ell\to k\) reduce through the same prime field.
An element of \(I_B\cap U\) supplied by \texttt{UFacts.mu}
also lies in \(\ker\eta\), so its image in
\(G_0=\ker\delta_\Phi\) proves \texttt{f2a}.
Push the \texttt{UFacts.f2b} witnesses to \(G_0\) to prove
\texttt{f2b}.  Push the \texttt{f2c} witnesses to \(\Gamma\);
their \(\delta_\Phi\)-value is one, so the additional
\(\delta_{\Phi,k}\) factor in the source is also one.
Every local witness in \texttt{UFacts.f3} lifts uniquely into
the indicated local decomposition group, proving \texttt{f3}.
An identity \(\operatorname{om}_\Phi=\psi_\Phi^j\) on
\(G_0\) would pull back to the identity forbidden by
\texttt{UFacts.f5}; this proves \texttt{f5}.

The source's
\(N_v=I_v\cap\ker(\epsilon_\Phi\circ\operatorname{loc}_v)\)
has \(\omega=1\).  On inertia above \(\ell\), the original
relation \(\alpha=\omega^j\) then makes \(\alpha=1\).
Equation (MFC.22) proves exactly \texttt{Psi\_N}.
These are all the asserted basic fields, including the actual
marked element and local residual character.

By MFC7 and \texttt{UFacts.deg}, \(d_v\ge\ell^m\).
There is at least one prime above \(\ell\), hence
\(d\ge\ell^m\).  MFC6's strict inequality now gives both
inequalities (MFC.23), since \(3C_0+7\ge7>5\).
\end{proof}
% END MFC9

\subsection{Two strict classes on this same quotient}

% BEGIN MFC10
\paragraph{MFC10. A nonzero residual Selmer class with one allowed prime.}
\label{claim:MFC10}
Assume the finite local and global laws and the Greenberg--Wiles
identity contained in \(P:\texttt{GC.Hypotheses}_{2}\).
Let \(\kappa:G\to k^\times\) have open kernel,
\(\kappa(\sigma)\ne1\), and \(\kappa\ne\omega\).
Let \(S\) be a finite set of rational primes containing
\(\ell\) and every ramified prime of \(\kappa\).
Suppose \(w\in S\), \(w\ne\ell\),
\(\kappa|_{D_A}=\omega|_{D_A}=1\) for \(A\mid w\),
and \(\kappa|_{D_A}\ne1\) at every other place in \(S\).
Then there is a locally constant cocycle
\[
 y(st)=y(s)+\kappa(s)y(t)                                  \tag{MFC.24}
\]
which is not a global coboundary, vanishes on inertia outside
\(S\), and is a coboundary on every decomposition group above
\(S\setminus\{w\}\).  This gives the full
\texttt{st\_sel} conclusion, including every prolongation.

\begin{proof}
Descend \(\kappa\) to \(G_{\mathbb Q,S}\): its kernel
contains each excluded inertia group and hence their closed
normal closure.  Regard \(k(\kappa)\) as a finite
\(\mathbb F_\ell\)-module.  Its global invariants vanish
because \(\kappa(\sigma)\ne1\).  Its cyclotomic twisted
dual has no global invariants either.  Indeed an invariant
\(\mathbb F_\ell\)-linear functional \(f:k\to\mathbb F_\ell\)
would satisfy \(f(\kappa(s)x)=\omega(s)f(x)\).
Choose \(s\) with \(\kappa(s)\ne\omega(s)\).  Multiplication
by their nonzero difference is bijective on \(k\), forcing
\(f=0\).  The real-place invariant group also vanishes.
At every \(v\in S\setminus\{w\}\), the assumed local
nontriviality gives zero local invariants by the same scalar
argument.

Impose local condition zero at those places and the full
\(H^1\) at \(w\).  Let \(\mathcal S\) and
\(\mathcal S^\perp\) be the resulting finite Selmer groups
and their dual conditions.  At \(w\), both the representation
and the residual cyclotomic character are trivial.  Local Tate
duality gives \(|H^2(\mathbb Q_w,k)|=|k|\), while
\(|H^0(\mathbb Q_w,k)|=|k|\).  The local Euler characteristic
away from \(\ell\) gives
\(|H^1(\mathbb Q_w,k)|=|k|^2\).
In the Greenberg--Wiles identity every global and real invariant
factor is one, and every local factor other than \(w\) is
one.  Its remaining equality is therefore
\[
                      |\mathcal S|=|k|\,|\mathcal S^\perp|.
                                                               \tag{MFC.25}
\]
The dual group contains zero, so its cardinality is at least
one, and \(|k|>1\).  Thus \(\mathcal S\) has a nonzero
class.  Choose a continuous representative and inflate it to
\(G\), obtaining \(y\).  Surjectivity of
\(G\to G_{\mathbb Q,S}\) ensures that this inflated cocycle
cannot be a global coboundary.  It is zero on excluded inertia,
and membership in the zero local conditions supplies the required
local primitives.

These primitives hold for every prolongation.  If
\(A'=\gamma A\) and \(y(s)=(\kappa(s)-1)b\) on \(D_A\),
the cocycle identity gives
\[
 y(\gamma s\gamma^{-1})
       =(\kappa(s)-1)(\kappa(\gamma)b-y(\gamma)).
\]
Thus the primitive at \(A'\) is
\(\kappa(\gamma)b-y(\gamma)\).  This verifies the exact
valuation scope, without a new independent local choice.
\end{proof}
% END MFC10

% BEGIN MFC11
\paragraph{MFC11. Restriction preserves nontriviality, and the auxiliary inertia is killed.}
\label{claim:MFC11}
Two elementary facts will be used.  First, if a subgroup
\(U\subseteq G\) contains \([G,G]\) and
\(\kappa|_U\ne1\), a \(\kappa\)-cocycle which is a
coboundary on \(U\) is already a coboundary on \(G\).
Second, at a rational prime \(p\ne\ell\), if an open
normal subgroup \(H\) has index \(\ell\) and
\(I_A\not\subseteq H\), then every locally constant additive
homomorphism \(y:I_A\to k\) vanishes on \(I_A\cap H\).

\begin{proof}
For the first assertion, subtract the supplied coboundary and
assume \(y|_U=0\).  The subgroup \(U\) is normal because it
contains the commutator.  Choose \(u\in U\) with
\(\kappa(u)\ne1\).  Comparing the two expressions
\(su=(sus^{-1})s\) in the cocycle identity gives
\(y(s)=\kappa(u)y(s)\).  The nonzero scalar
\(\kappa(u)-1\) is invertible, so \(y(s)=0\) for every
\(s\), as required.

For the second assertion, the quotient map
\(G\to G/H\simeq\mathbb F_\ell\) restricts to a nonzero
additive character \(\varphi\) of \(I_A\), with kernel
\(I_A\cap H\).  The maximal pro-\(\ell\) abelian quotient
of inertia at \(p\ne\ell\) is \(\mathbb Z_\ell\): the
wild inertia is pro-\(p\), and the tame inertia quotient has
one \(\mathbb Z_\ell\) factor.  Hence its continuous
\(\mathbb F_\ell\)-valued character space has dimension one.
Every coordinate of the additive map \(y\), in an
\(\mathbb F_\ell\)-basis of \(k\), is a multiple of
\(\varphi\).  It therefore vanishes on the same kernel.
All coordinates vanish there, giving the assertion for \(k\).
This is the source's \texttt{st\_tame} mechanism.  A trivial
valuation has trivial inertia and causes no exception.
\end{proof}
% END MFC11

% BEGIN MFC12
\paragraph{MFC12. Both complete \texttt{StrictCl} suppliers on one \(U_m\).}
\label{claim:MFC12}
For each \(a\in\{1,2\}\), put
\(\kappa_a=\alpha^{-a}\eta_k\).  Under \(P\), there is
a locally constant \(\kappa_a\)-cocycle \(y_a:G\to k\)
such that
\begin{align}
 &y_a|_{U_m}\text{ is not a coboundary};\notag\\
 &\forall A\mid\ell\ \exists b_A\quad
       y_a(s)=(\kappa_a(s)-1)b_A\quad(s\in D_A);\notag\\
 &A\nmid\ell,\quad s\in I_A\cap U_m,\quad
       \eta(s)=\alpha(s)=1\quad\Longrightarrow\quad y_a(s)=0.
                                                               \tag{MFC.26}
\end{align}
This is the full original \texttt{strict} theorem for each
\(m\) and each of its two allowed characters.  The two choices
are made on the same \(H_0,w,U_m\).

\begin{proof}
The character \(\kappa_a\) factors through \(\mu_{M^b}\)
for some \(b\), so its kernel is open and it is unramified
away from \(M\).  It is odd:
\(\kappa_a(\sigma)=-1\), since \(\alpha(\sigma)=1\).
It differs from \(\omega\): otherwise, on
\(U_m\cap\ker\eta\), one would have
\(\omega\alpha^a=1\), contradicting MFC4's appropriate
\texttt{f2b} or \texttt{f2c} witness.

At every \(A\mid\ell\), MFC4 gives \(s\in D_A\cap U_m\)
with \(\alpha(s)^2\ne1\).  Its \(\eta\)-value is one,
so \(\kappa_a(s)=\alpha(s)^{-a}\ne1\) for both \(a=1,2\).
At \(w\), the entire decomposition group kills
\(\eta,\alpha,\omega\) by (MFC.4).  Apply MFC10 with
\(S=\{\ell,w\}\) together with the ramified primes of
\(\kappa_a\).  At a ramified prime, the character is
nontrivial on inertia, and therefore on decomposition.  All
the hypotheses of MFC10 hold.  It yields a global noncoboundary
\(y_a\) with the indicated local primitives above \(\ell\).

At any prime other than \(w\), if
\(s\in I_A\) and \(\kappa_a(s)=1\), the cocycle is zero:
outside \(S\) it vanishes on inertia, while in
\(S\setminus\{w\}\) it is a local coboundary.
At \(w\), \(\kappa_a\) is trivial on decomposition, so
\(y_a|_{I_A}\) is additive.  The subgroup \(H_\theta(w)\)
has index \(\ell\), and inertia at \(w\) maps onto its
quotient, by CAF4.  MFC11 makes \(y_a\) vanish on
\(I_A\cap H_\theta(w)\), hence on \(I_A\cap U_m\).
This proves the last clause of (MFC.26) also at \(w\),
because \(\eta=\alpha=1\) implies \(\kappa_a=1\).

Finally \(\kappa_a\) is nontrivial on \(U_m\), witnessed
by the same local element used above.  Since \(U_m\) contains
the commutator, MFC11 says that a coboundary restriction to
\(U_m\) would make \(y_a\) a global coboundary.  This is
impossible.  All assertions of \texttt{StrictCl} are now verified.
Nothing in choosing \(y_1,y_2\) changes the already chosen
subgroup, field, local groups or marked element.
\end{proof}
% END MFC12

% BEGIN MFC13
\paragraph{MFC13. Exact descent of each strict class to the frame cone.}
\label{claim:MFC13}
Let \(\Phi\) be any actual linked frame as in MFC8, and let
\(y_a\) satisfy \texttt{StrictCl} for \(U=\operatorname{upl}(\Omega,L)\).
Then
\[
 H^1_C\!\left(\Phi.\operatorname{loc},
    k(\psi_\Phi^{-a}\delta_{\Phi,k}),\,b_v=0\right)
                         \ne0\qquad(a=1,2).               \tag{MFC.27}
\]
These are exactly the two source \texttt{Ok} fields
\texttt{f4} and \texttt{f4'}; the quotient and all local
maps are the same ones used in MFC8.

\begin{proof}
Write \(\chi_a=\psi_\Phi^{-a}\delta_{\Phi,k}\).
Equation (MFC.22) gives \(\chi_a\circ p=\kappa_a|_U\).
The strict inertia clause makes \(y_a\) zero on every
generator of the exact kernel \(J\) in (MFC.19).
To justify passage to its closed normal closure, consider
\[
                    H=\{u\in U:p(u)=1,\ y_a(u)=0\}.
\]
The cocycle identity makes \(H\) a subgroup.  If \(h\in H\),
then \(\kappa_a(h)=1\), and the same identity gives
\(y_a(uhu^{-1})=\kappa_a(u)y_a(h)=0\); hence it is normal.
It is closed because \(p\) is continuous and \(y_a\) is
locally constant.  It contains the kernel generators, so it
contains \(J\).  Thus \(y_a\) is constant on fibers of
\(p\), and defines \(\varphi_a:\Gamma\to k\).
The quotient-map property of \(p\) shows that
\(\varphi_a\) is locally constant.  Its cocycle identity
follows by lifting two elements of \(\Gamma\) to \(U\).

For each \(v\), the chosen \(A_v\) lies above \(\ell\).
Choose its local primitive \(b_v\) from (MFC.26).  Under
the exact inclusion \(j_v\), this gives
\[
 \varphi_a(\operatorname{loc}_v(s))
      =(\chi_a(\operatorname{loc}_v(s))-1)b_v
                                  \quad(s\in G_v).        \tag{MFC.28}
\]
By the literal source definition, the pair
\((\varphi_a,(b_v)_v)\) is therefore a member of
\(Z^1_C\) with all \(b_v\)-subspaces equal to zero.
If its class in \(H^1_C\) were zero, the source's
\texttt{continuousCoboundaries1} condition would supply one
\(b\in k\) with
\(\varphi_a(g)=(\chi_a(g)-1)b\) for every \(g\).
Pulling back by \(p\) would make \(y_a|_U\) a coboundary,
contrary to MFC12.  Thus the class is nonzero, proving (MFC.27).
The argument applies to both characters without changing any
frame data.
\end{proof}
% END MFC13

% BEGIN MFC14
\paragraph{MFC14. The actual shared marked construction, with an exact remaining boundary.}
\label{claim:MFC14}
Fix the full original \texttt{Gen}, the original \(P\), and
any prescribed \(Q>0\).  There exist one
\[
                  L,\quad g_1,\quad\Phi,\quad C_0,\quad w
\]
with the complete original \texttt{Linked} relation,
\texttt{Big}\(\ C_0\), the witness-prime conditions
(MFC.4), \(w\equiv1\pmod Q\),
\(\operatorname{upl}(\Omega,L)\subseteq H_\theta(w)\),
and all the twenty-one \texttt{Ok} fields proved in MFC9
and MFC13.  The bound (MFC.13) holds at the selected layer.
When \(O=c.O\) is the original coefficient DVR, choose
\(Q\) first by CAF7.  This same constructed \(L\) then
has the full original \texttt{SeedClause}, with the same
marked square witness \(g_1\).

\begin{proof}
Choose \(H_0\) by MFC1 and choose \(w\) by CAF5 for the
already fixed \(Q\).  Choose \(C_0\) by MFC5, before
selecting \(m=3C_0+7\).  MFC4 and MFC6 give the actual
embedded CM field and marked \(g_1\).  MFC8 constructs
\(\Phi\) with every linking field.  MFC9 proves the basic
fields and \texttt{Big}; MFC12 and MFC13 prove the two
remaining strict-class fields on this very quotient.
The subgroup inclusion follows from the definition of \(U_m\),
and the class-rank bound is the selected instance of MFC5.

For the coefficient-DVR assertion, the full \texttt{Gen}
implies the seed input used by CAF7.  Make its choice of
\(Q\) before the preceding construction.  The resulting
\(L,w\) satisfy all its field, splitting and congruence
hypotheses.  Equation (MFC.16) supplies
\(g_1\in G_L\) with \(\alpha(g_1)^2\ne1\), so CAF7
returns the original \texttt{SeedClause}.  Its character
triple is installed on this field; it does not replace the
field, marked element, quotient, local maps or either strict class.
This completes the claimed construction at its stated boundary.
\end{proof}
% END MFC14

\subsection{The seven remaining \texttt{Ok} fields, and the original exits}

The seven fields still needed on the \emph{same} \(\Phi,C_0\)
are \texttt{phi}, \texttt{gen}, \texttt{cl}, \texttt{ec},
\texttt{b1p}, \texttt{b1} and \texttt{h1}.
Their source scopes remain unchanged: \texttt{phi} is the
global \texttt{PhiCond}; \texttt{gen} chooses \(d_v\) local
inertia generators with the prescribed exceptional index and
character restrictions, uniformly for every finite abelian
\(\ell\)-group character; \texttt{cl} gives one exponent
annihilating every everywhere-unramified finite abelian
\(\ell\)-group character of \(\Phi.G0\); and \texttt{ec}
is the full cone Euler inequality for every finite coefficient
field, finite representation and stable collection of local
subspaces satisfying its stated dual-invariant condition.

The last three retain their exact coefficient modules and local
conditions:
\begin{align}
 \operatorname{rank}_k H^1_C(\operatorname{loc},
       k(\psi_\Phi^{-2}\delta_{\Phi,k}),0)&\le C_0,
          &&\texttt{b1p},\notag\\
 \operatorname{rank}_k H^1_C(\operatorname{loc2},
       k(\psi_\Phi^{-1}|_{\Phi.G0}),0)&\le C_0,
          &&\texttt{b1},\notag\\
 \operatorname{rank}_k H^1_C(\operatorname{loc},k,\top)
                                                &\le1+C_0,
          &&\texttt{h1}.                                  \tag{MFC.29}
\end{align}
The bare uniform unramified character bound in MFC5 is not
silently identified with these twisted cone bounds.  Their
restriction, coefficient and local-condition comparisons still
have to be supplied.  Filling these seven fields in the same
constructed record gives the literal \texttt{HasFrameW};
MFC14's seed installation then gives the literal
\texttt{HasFrame}.

To obtain the original \texttt{GdFrameEx} and
\texttt{HostFrameEx}, their full entry suppliers
\texttt{gen\_gd} and \texttt{gen\_host} must also be connected
at their original scopes.  CAF9--CAF10 supply their seed subset,
not every \texttt{Gen} field.  In particular the host's
\(H_b=\ker\kappa_B\), its index, local lifting property
and nonquadratic decomposition witness must be retained; a
\(\top\) subgroup or the Dat numerical restriction cannot
be substituted.  The rank-two/rank-three automorphic seeds,
common coefficient extension and original/alternative endpoint
matching remain further work beyond these arithmetic frames.

\paragraph{Pinned correspondence.}
The immutable primary sources are the \texttt{FrobeniusSmoothVectors}
files named in this chapter, the literal frame and linking definitions
in \texttt{Continuous\allowbreak Coboundaries2}, \texttt{HostFrameEx} and
\texttt{GLinkQuot}, and the exact class-character definition
\texttt{Kum.XDual}.  Complete file snapshots, statement ranges,
source hashes and claim-block hashes accompany this chapter.
The elementary CRT choice in MFC1 is an intentional human
simplification with the same complete output record.  No changed
field model, local identification or order of choices is hidden
in that simplification.
